%% REV_C2: section file; no preamble or document environment.
\section{Finite trail and stationary Bellman certificates}

Throughout this appendix the thresholds are
\(c_{\rm s}=961/1000\) and \(c_{\rm t}=187/200\), so that
\(x_{\rm s}=1000/961\) and \(x_{\rm t}=200/187\).
The trail weights, response pairs, orderings and vector \(h\) are
those of Definitions~\ref{d:weights}, \ref{d:kernel} and
\ref{d:kernel-order}. In particular, a revisit means an earlier
\emph{departure} from the same site. All positions below are points
of \(\mathbb Z^2\).

The finite assertions in Propositions~\ref{p:H2-test},
\ref{p:automaton10} and \ref{p:stationary-check} specify exact checks
of the archived witnesses. The arguments below prove the soundness
of those checks. The finite checks themselves are stated with their
full arithmetic contracts, for separate verification; an executable's
exit code or a floating approximation is not an additional hypothesis
of any soundness lemma.

\begin{definition}\label{d:memory-automaton}
Let \(\mathcal E=\{{\rm E}=(1,0),{\rm N}=(0,1),
{\rm W}=(-1,0),{\rm S}=(0,-1)\}\).
For \(m\ge1\), \(\Omega_m\) consists of sequences
\(\omega=(d_1,\ldots,d_m)\in\mathcal E^m\) whose path from zero
uses distinct unoriented bonds. Write \(z_0=0\),
\(z_j=\sum_{i\le j}d_i\), and \(a(\omega)=d_m\).
A direction \(d\) is admissible if \(d\ne-d_m\) and
\(\{z_m,z_m+d\}\) is not a used bond of this window.
Its successor is \(\omega d=(d_2,\ldots,d_m,d)\).
Its kind \(k(\omega,d)\) is \({\rm r}\) if
\(z_m\in\{z_0,\ldots,z_{m-1}\}\), otherwise \({\rm s}\)
if \(d=d_m\), and \({\rm t}\) otherwise.
Direction labels remain absolute lattice directions.
\end{definition}

The successor belongs to \(\Omega_m\): the last \(m-1\) old bonds
and the new bond are a subset of the distinct bonds of the extended
window, and translation to a new initial point preserves bond equality.
The explicit nonreversal condition is redundant with the used-bond
test, since the last bond is used, but fixes the incoming-direction
convention.

\begin{definition}\label{d:memorygrowth}
On \(\Omega_m\), define the nonnegative one-step matrix
\[
 \mathsf A_m(\omega,\omega')=
 \sum_{\substack{d\ {\rm admissible}\\\omega d=\omega'}}
 \gamma(\omega,d),\qquad
 \gamma(\omega,d)=
 \begin{cases}1/5,&k={\rm r},\\
 x_{\rm s}/3,&k={\rm s},\\x_{\rm t}/3,&k={\rm t}.
 \end{cases}
\]
For a prescribed start and incoming direction with initial reversal
forbidden, let \(A_n(\omega)\) be the true total \(\mathsf W\)-weight
of length-\(n\) trails whose last \(m\) steps form \(\omega\),
for \(n\ge m\), and let \(W_n=\sum_{|w|=n}\mathsf W(w)\).
The adopted finite test is \(m=6\) and extension length \(1\).
\end{definition}

There is no virtual incoming bond in the used-bond set. Only the
initial reversal is forbidden. Thus the first departure has the same
straight or turning convention as every later first departure.

\begin{lemma}\label{l:alm}
For \(n\ge m\),
\[
 A_{n+1}(\omega')\le\sum_\omega A_n(\omega)\mathsf A_m(\omega,\omega').
\]
If \(v_\omega>0\), \(0<q<1\), and \(\mathsf A_mv\le qv\),
then (H2) holds with
\[
 C_V=\max\left\{
      \max_{0\le r<m}W_rq^{-r},\
      \frac{\max v}{\min v}W_mq^{-m}\right\}.
\]
Any larger finite constant is also valid.
\end{lemma}

\begin{proof}
Fix the initial data and a true trail \(w\) of length \(n+1\).
Its length-\(n\) prefix is a true trail with the same initial data.
If its last window is \(\omega\), its next direction is admissible
for \(\omega\): a used bond of this window is a used bond of the
true trail, and reversal would reuse the last bond. The last window
of the extension is exactly \(\omega d\). The map from \(w\) to its
prefix and final direction is injective.

Suppose first that the window detects a revisit. Its current endpoint
then equals an earlier window vertex from which a window bond departs.
This is also an earlier true departure. The new true and window
weights are both \(1/5\). If the true departure is a revisit not
detected by the window, its weight is \(1/5\), whereas the window
uses \(x_{\rm s}/3\) or \(x_{\rm t}/3\), each at least \(1/3\).
Otherwise both departures are first departures; their incoming
direction is the last window direction, so their type and weight
agree. These cases exhaust all departures. Multiplication by the
nonnegative prefix weight and summation over the injective map give
the asserted componentwise inequality.

Write \(A_n\) as a row vector and put
\(v_- =\min_\omega v_\omega>0\),
\(v_+=\max_\omega v_\omega\). Nonnegativity of the matrix implies,
by induction on \(j\),
\[
 A_{m+j}\le A_m\mathsf A_m^j,\qquad
 \mathsf A_m^jv\le q^jv .
\]
Consequently
\[
 W_{m+j}\le v_-^{-1}A_m\mathsf A_m^jv
       \le (v_+/v_-)q^jW_m .
\]
This proves the claimed bound for all lengths at least \(m\).
For lengths less than \(m\) it is the defining finite maximum.
The empty trail contributes \(W_0=1\), so \(C_V\ge1\).
Translations and the eight symmetries of the square give
weight-preserving bijections between the trail sets for any two
prescribed starts and incoming directions. Thus all the base sums
\(W_r\) in this formula are uniform over those initial data.
Increasing \(C_V\) preserves the inequality, which proves
\ref{h:H2} with the stated constants.
\end{proof}

\begin{lemma}\label{l:H2-dihedral}
The dihedral group of the square acts on \(\Omega_m\) by applying
the same lattice symmetry to every direction. This action preserves
admissibility, kind, successor and the weights of \(\mathsf A_m\);
thus \(\mathsf A_m(g\omega,g\omega')=\mathsf A_m(\omega,\omega')\).
For any positive vector constant on these orbits, a row inequality
\((\mathsf A_mv)_\omega\le qv_\omega\) checked at one representative
of every orbit holds on every raw window.
For \(m=6\), \(|\Omega_6|=916\) and there are \(115\) orbits.
\end{lemma}

\begin{proof}
A square symmetry is an invertible linear map carrying the four
unit directions to themselves. It carries vertices and unoriented
bonds bijectively to vertices and bonds. It therefore preserves
bond distinctness, occurrence of the endpoint among earlier
departure vertices, equality of the outgoing and incoming directions,
and the reversal relation. It also commutes with deleting the first
direction and adjoining the last. This proves every asserted
transition property, and hence the matrix identity, by the bijection
\(d\mapsto gd\) in its defining sum. If \(v\) is constant on orbits,
then
\[
 (\mathsf A_mv)_{g\omega}
 =\sum_{\omega'}\mathsf A_m(g\omega,g\omega')v_{g\omega'}
 =(\mathsf A_mv)_\omega ,\qquad v_{g\omega}=v_\omega .
\]
This proves the row reduction.

Here is a direct count for memory six. Let \(N_j=|\Omega_j|\),
with \(N_0=1\). There are \(N_1=4\) initial directions and
three nonreversing choices at each later step before any additional
bond obstruction. A nonreversing square-lattice path cannot reuse a
bond before step five: such reuse would require a return to a previous
vertex, and the shortest nonreversing cycle has length four.
Thus \(N_2=12,N_3=36,N_4=108\).
At step five the only additional obstruction is a four-step square
starting at the initial vertex, followed by its first edge. There
are four choices of first direction and two orientations of that
square. Exactly these eight extensions must be deleted, giving
\(N_5=3\cdot108-8=316\).

For a valid length-five prefix, an additionally forbidden sixth
step must start at its vertex of time one. Indeed, a repeat of an
earlier vertex at time five has an even time difference at least
four, and only time one qualifies; a difference two would already
reverse a bond. Hence steps two through five form a square based at
that vertex. There are initially \(4\cdot3\cdot2=24\) choices
of the first direction, the square's first direction and its
orientation. Of these, exactly eight have the square's last edge
equal to the first bond of the entire path: for each first
direction, the square's first direction must be one of its two
perpendicular directions, and precisely one orientation then ends
along that first bond. These eight are invalid length-five
prefixes. Each of the remaining sixteen prefixes has precisely two
used edges at the current vertex besides the reverse of the most
recent edge: the initial bond and the first edge of the square.
Both give distinct forbidden sixth directions. This gives
\(N_6=3\cdot316-2\cdot16=916\).

For completeness, count the orbits without a supplied representative
list. The number of orbits is
\(\frac18\sum_g|\{\omega:g\omega=\omega\}|\): to see this,
sum the number of stabilizers over each orbit and use that the orbit
size times its stabilizer size is eight. A nonidentity rotation, or
a reflection in a diagonal, fixes no unit direction and hence no
window. Reflection in either coordinate axis fixes two directions.
A window using just those two opposite directions must use a single
one throughout, since a change of sign immediately reverses a bond.
Each such reflection therefore fixes exactly two windows. The orbit
count is \((916+2+2)/8=115\).
\end{proof}

\begin{proposition}\label{p:H2-test}
At the two fixed H1 thresholds there exists a positive integer
vector \(v\) on \(\Omega_6\), constant on dihedral orbits, with
\[
 \frac{\max v}{\min v}=\frac{250000000000}{46594269069}<6,\qquad
 \max_\omega\frac{(\mathsf A_6v)_\omega}{v_\omega}
       =q_*:=\frac{1119948287353}{1120437614225}<\frac{2499}{2500}.
\]
The weights have common denominator \(D_2=2695605\), with
numerators \(935000,961000,539121\) for \({\rm s},{\rm t},{\rm r}\).
For \(q_*=q_{\rm num}/q_{\rm den}\), the required row test is
\[
 q_{\rm den}\sum_{d\ {\rm admissible}}n_{k(\omega,d)}v_{\omega d}
          \le q_{\rm num}D_2v_\omega .
\]
It suffices to check the \(115\) orbit rows after proving
Lemma~\ref{l:H2-dihedral}; checking all \(916\) raw rows is also valid.
\end{proposition}

\begin{proof}[Finite verification contract]
Number the directions E,N,W,S by \(0,1,2,3\), and for this
certificate use the least-significant-first code
\[
 e_6(d_1,\ldots,d_6)=\sum_{i=1}^6d_i4^{i-1}.
\]
The canonical code is the minimum of \(e_6\) over the eight images
\(d_i\mapsto \varepsilon d_i+r\pmod4\),
\(\varepsilon\in\{1,-1\}\), \(r\in\{0,1,2,3\}\).
Use the integer field \texttt{h} indexed by the field \texttt{code}
in \texttt{h2\_exact\_m6k1\_r5.json}; pull it back by this
canonical-code map to obtain \(v\). The decimal vector files are
not used to round or reconstruct these integers.

Generate windows by starting with the empty word at zero and,
at each stage, adjoining every direction whose unoriented bond has
not yet been used. Each generated word is a trail. Conversely,
deleting the last step of a trail gives a previously generated word,
so induction on its length proves that every window is generated
exactly once. Lemma~\ref{l:H2-dihedral} proves that the canonical
codes give exactly the needed orbits. The finite data checks are:
the stored codes are precisely those \(115\) canonical codes,
without repetition; their integer values are positive; and every
regenerated row satisfies the displayed integer inequality.
In the three cases the weights are exactly
\[
 \frac{1000}{2883}=\frac{935000}{2695605},\qquad
 \frac{200}{561}=\frac{961000}{2695605},\qquad
 \frac15=\frac{539121}{2695605}.
\]
Thus the checked sums are the rows of the mathematically defined
matrix, with no matrix supplied as an input.

To specify also the equalities in the proposition, the integer
extremum checks are
\[
 \min v=186377076276\quad(\text{canonical code }441),\qquad
 \max v=1000000000000\quad(\text{canonical code }0).
\]
At canonical code \(428\), the regenerated row has
\[
 v_{428}=224087522845,\qquad
 \sum_d n_{k(\omega,d)}v_{\omega d}=603787640626036713.
\]
Exact cancellation gives
\[
 \frac{603787640626036713}{2695605\cdot224087522845}
 =\frac{1119948287353}{1120437614225}=q_*.
\]
The upper row checks and this equality establish the asserted
maximum; the two extremum checks establish its asserted vector
ratio. The remaining strict comparisons follow by multiplying
positive integer denominators. These are finite integer checks on
the archived witness, rather than assumptions about a computed
eigenvalue or a numerical power iteration.
\end{proof}

\begin{lemma}\label{l:H2-CV}
With \(m=6\), \(q=q_*\), and the vector in Proposition~\ref{p:H2-test},
the constant supplied by Lemma~\ref{l:alm} is at most \(10\).
Indeed, uniformly over the prescribed start and incoming direction,
\[
 W_r\le x_{\rm t}^r\quad(0\le r\le6),\qquad
 \max\left\{\max_{r<6}W_rq^{-r},\
                 \frac{\max v}{\min v}W_6q^{-6}\right\}
 \le6\left(\frac{200}{187q}\right)^6<10.
\]
Consequently (H2) holds with the explicit rational constants
\(q=q_*\) and \(C_V=10\).
\end{lemma}

\begin{proof}
The first step has at most three choices because the designated
initial reversal is forbidden. At every subsequent step the last
bond is used, again leaving at most three choices. Every allowed
weight is at most \(x_{\rm t}/3\), since
\(1/5\le1/3\le x_{\rm s}/3\le x_{\rm t}/3\).
For each prefix, its total extension weight is therefore at most
\(x_{\rm t}\) times its weight. Induction, starting with \(W_0=1\),
gives \(W_r\le x_{\rm t}^r\), in particular for \(0\le r\le6\).

Put \(b=x_{\rm t}/q>1\). For \(r<6\),
\(W_rq^{-r}\le b^r\le b^6\), and the second term in
Lemma~\ref{l:alm} is at most \(6b^6\) by
Proposition~\ref{p:H2-test}. This proves the non-strict bound
in the statement. To verify its last inequality by small integer
comparisons, cross multiplication gives
\(q_*>999/1000\), and hence
\[
 \frac{200}{187q_*}<\frac{200000}{186813}<\frac{15}{14}.
\]
The second inequality is \(2800000<2802195\).
Finally
\(6\cdot15^6=68343750<75295360=10\cdot14^6\).
Thus \(6b^6<10\), and Lemma~\ref{l:alm} permits \(C_V=10\).
\end{proof}

\begin{lemma}\label{l:automaton-quotient}
Let \(\pi:\Omega_m\to Q\) be a surjection onto a finite set.
Suppose that windows with the same image have, for each absolute
direction, identical admissibility and kind, and their admissible
successors have the same image. Suppose also that two such windows
have the same final direction whenever their class has an admissible
non-revisit transition.
Then admissibility, kind and successor define a well-defined automaton
on \(Q\). Every phase-and-position transition whose coefficients
depend only on the absolute direction, kind, position, and final
direction on non-revisit edges factors through this automaton.
For every function \(G(\varphi,q,x)\), its pullback by \(\pi\)
intertwines the raw and quotient backward operators. This
intertwining also holds for each finite stopped continuation sum.
\end{lemma}

\begin{proof}
For \(q\in Q\), choose any \(\omega\) with \(\pi(\omega)=q\).
Declare direction \(d\) admissible at \(q\) exactly when it is
admissible at \(\omega\), give it kind \(k(\omega,d)\), and give
it successor \(\pi(\omega d)\). The hypotheses make all three
choices independent of \(\omega\). When the coefficient of a
non-revisit edge needs the incoming direction, use \(a(\omega)\).
That direction is independent of the representative in every class
with such an edge. In a class with only revisit edges no coefficient
needs it. This proves well-definedness even when such a class
contains several final directions.

Here is the operator assertion including possible parallel edges.
At a raw state \((\varphi,\omega,x)\), write its backward row as
the finite sum over admissible directions \(d\) and allowed next
phases \(\psi\), with coefficient
\(t(\varphi,\omega,x;d,\psi)\). The quotient row has the same
summation indices, coefficient and new position \(x+d\), and its
new window is \(\pi(\omega d)\). If
\(\pi^*G(\varphi,\omega,x)=G(\varphi,\pi(\omega),x)\), termwise
comparison gives
\[
 T_{\rm raw}(\pi^*G)=\pi^*(T_QG).
\]
There is no division by a class size: this is an equality of
backward rows, retaining every direction separately.

For stopped sums the terminal reward and stopping rule must also
factor through the displayed state, as they do for all the terminal
sets used here. The zero-step reward then pulls back, and the
intertwining identity inductively gives
\(T_{\rm raw}^jg_{\rm raw}=\pi^*(T_Q^jg_Q)\) for each \(j\).
Summing finitely many such identities proves the stopped-sum
assertion. No property of a partition-minimization algorithm is used.
\end{proof}

\begin{proposition}\label{p:automaton10}
There are \(60092\) raw windows in \(\Omega_{10}\).
The supplied map to \(Q=\{0,\ldots,3740\}\), with its supplied
representative per class, is surjective and satisfies every
hypothesis of Lemma~\ref{l:automaton-quotient}.
Its admissible labelled edges consist of \(2516\) straight,
\(5032\) turning and \(1872\) revisit edges.
\end{proposition}

\begin{proof}[Finite reconstruction contract]
For this file, in contrast to the H2 orbit file, the direction code
is most-significant-first:
\[
 e_m(d_1,\ldots,d_m)=\sum_{i=1}^m d_i4^{m-i}.
\]
The header of \texttt{windows\_m10\_min.txt} is
\((60092,3741)\). It is followed by \(60092\) pairs
\((e_{10}(\omega),\pi(\omega))\), and then one row per class.
A class row consists of its representative code, four successor
class numbers and four kinds, in absolute order E,N,W,S.
An inadmissible direction has successor and kind both \(-1\);
the kind numbers \(0,1,2\) mean \({\rm s},{\rm t},{\rm r}\).

The generating recursion in the proof of
Proposition~\ref{p:H2-test}, now continued to depth ten and ordered
lexicographically, defines the index set independently of the file.
Its finite cardinality checks, including useful intermediate counts,
are
\[
 \begin{array}{c|rrrrrrrrrr}
 j&1&2&3&4&5&6&7&8&9&10\\ \hline
 |\Omega_j|&4&12&36&108&316&916&2628&7500&21268&60092.
 \end{array}
\]
As proved for that recursion, no word is missing or duplicated.
The remaining finite assertions are these: the map rows list exactly
these words in that order; all class numbers lie in
\(\{0,\ldots,3740\}\); each representative is a generated word
and maps to its own class; and the file has exactly the stated
number of rows.

For every raw word, regenerate its endpoint, its used bonds, and its
earlier departure vertices. For each of its four directions, compare
admissibility and kind with its representative, compare the image
of its successor with the stored successor, and, for every admitted
non-revisit direction, compare final directions. The finite check
requires equality in each comparison. Thus it verifies precisely
the hypotheses of Lemma~\ref{l:automaton-quotient}, including the
final-direction condition. Representatives mapping to their classes
prove surjectivity. Counting the admitted directions in the class
rows gives the integer assertions \(2516,5032,1872\), respectively.
This specifies the finite automaton check completely and makes no
minimality or spatial-symmetry assertion.
\end{proof}

\begin{definition}\label{d:phaseweights}
The phases are \({\rm A},{\rm B},{\rm C},{\rm D}\). Their
diagonal transition weights, in kind order \({\rm s},{\rm t},{\rm r}\), are
\[
 \begin{array}{c|ccc}
 {\rm A}&x_{\rm s}/3&x_{\rm t}/3&1/5\\
 {\rm B}&x_{\rm s}/3&1/3&1/5\\
 {\rm C}&x_{\rm s}/3&1/3&1/5\\
 {\rm D}&1/3&1/3&1/5
 \end{array}
\]
At a first visit with reward \(\kappa_h\), additionally allow
\({\rm A}\to{\rm C}\) with weight \(x_{\rm t}\kappa_h/3\) on a turn,
or \({\rm B}\to{\rm D}\) with weight \(x_{\rm s}\kappa_h/3\)
on a straight passage. No other phase change is allowed.
At a revisit there is no phase change.
Start with weight one in each of \({\rm A},{\rm B}\);
the terminal reward is \(g=(0,0,1,1)\).
For a true trail, use its actual revisit kinds and its actual
type-dependent live-pair reward from Definition~\ref{d:kernel}.
\end{definition}

Here \(\kappa_h=\sum_{Y,s}\kappa_{Y,s}h_Y\). In the free
three-neighbour rule its two possible values are
\[
 \kappa_{h,\rm s}^{\circ}=a_{\rm s}+2b_{\rm s}h_C,
 \qquad
 \kappa_{h,\rm t}^{\circ}=b_{\rm t}+(a_{\rm t}+b_{\rm t})h_C.
\]
Indeed, the pair of non-outgoing neighbours is opposite on a
straight passage and perpendicular on a turn; among the two pairs
with the outgoing neighbour, both are perpendicular in the straight
case, and one is opposite in the turning case. A sole surviving
pair instead gives \(b'_\tau h_Y\) with its actual pair type \(Y\).

\begin{lemma}\label{l:phase-expansion}
For a true closed trail \(w\) ordered with straight first visits
before turning first visits, its \(h\)-weighted contribution to
\(M_1\) equals the sum over phase histories starting in
\({\rm A}\) or \({\rm B}\) and ending in \({\rm C}\) or \({\rm D}\)
of the products of the weights in Definition~\ref{d:phaseweights}.
Equivalently this sum is
\[
 3^{-|w|}(3/5)^{\operatorname{rev}(w)}
       \sum_k\kappa_h(\sigma_w(k))\prod_{l\le k}x_{\sigma_w(l)}.
\]
Splitting such a phase history at any prefix length gives the
product of its exact prefix weight and its corresponding suffix weight.
\end{lemma}

\begin{proof}
Put \(n=|w|\). An allowed phase history which starts in A and
has nonzero terminal reward must change from A to C exactly once.
There is no other way to enter C, and no transition leaves C.
Its change occurs at a unique turning first departure \(j\).
Before that departure the diagonal factors are
\(x_{\rm s}/3,x_{\rm t}/3,1/5\) on straight, turning and revisit
departures. At \(j\) the factor is
\(x_{\rm t}\kappa_h(j)/3\). After it the factors are
\(x_{\rm s}/3,1/3,1/5\). Their product is therefore
\[
 3^{-n}(3/5)^{\operatorname{rev}(w)}\kappa_h(j)
 x_{\rm s}^{\#\{\text{straight first departures}\}}
 x_{\rm t}^{\#\{\text{turning first departures at indices }i\le j\}}.
\]
These are precisely the first departures at or before \(j\) in
the straight-first ordering, including all straight first departures
even if their path indices exceed \(j\).

A history starting in B with nonzero terminal reward changes from
B to D exactly once, at a straight first departure \(j\). Its
factors before, at and after that departure similarly give
\[
 3^{-n}(3/5)^{\operatorname{rev}(w)}\kappa_h(j)
 x_{\rm s}^{\#\{\text{straight first departures at indices }i\le j\}}.
\]
These are exactly the first departures at or before that \(j\) in
the same ordering. No revisit can be the changing departure.
There is thus one allowed history for each candidate responding
first departure; if its reward is zero both corresponding terms
are zero. Summing over all such departures gives the displayed
identity. For long trails it is the ordering fixed in
Definition~\ref{d:kernel-order}, hence the claimed contribution
to \(M_1h\). For other trails this identity concerns their stated
straight-first ordering.

Finally split a history after \(r\) departures, \(0\le r\le n\).
Its edge-weight product is the product of the factors with indices
less than \(r\) and those with indices at least \(r\).
Grouping by the phase at the split proves the factorization.
When used for a true suffix, its weights retain the prefix history
in deciding revisits and dead neighbours. The assertion requires
no independence or loss of that history.
\end{proof}

\begin{definition}\label{d:stationary-walk}
Fix \(X=B,C\), \(d_0\in\{{\rm N},{\rm W},{\rm S}\}\),
the sites \(u,c,v\) of Definition~\ref{d:kernel}, and \(a=u+d_0\).
The state space is
\(\{{\rm A},{\rm B},{\rm C},{\rm D}\}\times Q\times\mathbb Z^2\).
At positions \(c,v\) there are no outgoing transitions.
At any other position \(x\), an admissible automaton direction \(d\)
goes to position \(x+d\) and the quotient successor window.
At \(x=u\) or \(a\), forbid reuse of \(\{u,a\}\) in either
direction and kill every memory-detected revisit departure;
every other allowed departure there has weight \(1/5\),
in the same phase and with no response.
At other positions use the automaton kind and the phase weights.
For a non-revisit response, declare only \(c\) dead among the three
neighbours other than the arrival site, and use the type-dependent
three-neighbour or two-neighbour reward rule of Definition~\ref{d:kernel}.
A transition into \(c\) has zero contribution.
Let \(T_{X,d_0}\) be this nonnegative backward operator, with zero
terminal rows, and \(g(\varphi,q,x)=\mathbf1_{\{x=v\}}g_\varphi\).
The free operator \(\mathsf W\) omits positions and special-site
rules, uses the same quotient automaton and phase weights, and
treats all three response neighbours as live.
\end{definition}

Directions remain separate summation indices even if two automaton
edges share a successor class. The position displacement of an edge
is its actual direction. Thus, for a function \(H\), the operator
means the finite sum
\[
 (T_{X,d_0}H)(\varphi,q,x)
   =\sum_{d,\psi}t_{X,d_0}(\varphi,q,x;d,\psi)
                         H(\psi,q d,x+d),
\]
where forbidden transitions have coefficient zero. One may delete
edges into the centre because its continuation value is zero.
The quotient operator is well-defined by
Lemma~\ref{l:automaton-quotient}.

\begin{lemma}\label{l:stationary-domination}
Let \(P=(v_0,\ldots,v_L)\) be a true prefix avoiding \(c,v\),
with \(v_0=u\), \(v_1=u+d_0\), and \(L\ge m+2\), \(m=10\).
Every true continuation to its first hit of \(v\), together with
each phase history, is dominated by the corresponding path of
\(T_{X,d_0}\) from its prefix phase, last-window class and endpoint.
The comparison holds for every finite suffix horizon and hence
for their nonnegative suprema.
In this comparison:
forgotten true revisits have weight \(1/5\) no greater than any
relaxed first-visit diagonal weight; forgetting dead non-outgoing
neighbours only increases the reward, using \(b'_\tau\le b_\tau\);
and the known-site killing rules cannot kill a true continuation.
Indeed, the known initial departures at times \(0,1\) are outside
the memory at every suffix time, and a memory-detected additional
departure at \(u\) or \(u+d_0\) would be a forbidden third departure.
\end{lemma}

\begin{proof}
Fix a true completed trail extending \(P\), of length \(n>L\),
and one of its phase histories. At a suffix departure time
\(j\), \(L\le j<n\), its last \(m\) directions form a raw
window. The next bond is not a used window bond, and reversing
would reuse its last bond. Its direction is therefore admitted by
the raw automaton and, by Lemma~\ref{l:automaton-quotient}, by the
quotient. The new window class is the quotient successor.
The centre is never visited. The target is not visited before time
\(n\); its terminal reward is exactly the phase reward. Thus the
centre and target stopping rules discard none of these histories.

We verify the permanent initial-site rules explicitly. The bond
\(\{u,a\}\) was used at time zero and cannot be reused in either
direction by a true trail. Both \(u\) and \(a\) have already
been departed from, at times zero and one respectively; these
departures exist because \(L\ge12\), and neither site is the
target. Any suffix departure from either site is consequently a
true revisit and has diagonal weight \(1/5\) and no response.
At a suffix time \(j\), the remembered departure indices are
\(j-m,\ldots,j-1\). Since \(j-m\ge L-m\ge2\), these do not
include times zero or one. A detected revisit at one of the two
known sites would therefore certify another departure there in
addition to its initial departure and its current departure.

Such a third departure is impossible by distinctness of the used
bonds. At \(u\), the first departure consumes one bond, and each
later departure requires a distinct arrival bond and a distinct
departure bond. Three departures would consume at least five
bonds; only three are available when \(c\) is excluded.
Here \(a\) is one of \((-1,1),(-2,0),(-1,-1)\), so it is distinct
from \(u,c,v\) in both rows. At \(a\), the initial arrival and
departure consume two
bonds, and two further departures would each consume two additional
bonds, requiring six among four. Consequently the detected-revisit
killing rule cannot kill the chosen history. An allowed departure
from a known site has exactly its true weight \(1/5\).

Now suppose the departure site is neither known site. If the
window detects a revisit, it contains an earlier departure from
that site; this is a true revisit, so the diagonal factors agree
and no phase change occurs. If the true departure is a revisit
which the window has forgotten, keep the same phase in the relaxed
history. Every relaxed first-visit diagonal factor is at least
\(1/3\), and hence at least the true \(1/5\). No responding
departure is lost in this case, because a true revisit never responds.

At a true first departure the window also has a first departure,
with the same incoming and outgoing directions. Its diagonal factors
agree. For a responding departure compare the rewards as follows.
The centre is truly dead, so every neighbour declared dead in the
relaxation is truly dead. If both true non-outgoing neighbours are
live, the relaxed and true pair rules coincide. If precisely one is
live and the relaxation retains the other as dead, the rules again
coincide. If the relaxation forgets that dead neighbour, its
coefficient on the existing pair increases from \(b'_\tau\) to
\(b_\tau\), and its other added pair coefficients are nonnegative.
If both true non-outgoing neighbours are dead, the true reward is
zero. These cases, together with \(0<b'_\tau\le b_\tau\) and
\(h_B,h_C>0\), prove that the relaxed response factor is at least
the true response factor. The outgoing neighbour cannot be dead
on a true continuation, and a true transition never enters \(c\).

Multiplying the comparisons proves domination of the chosen
history. For a fixed prefix and a specified phase at the split,
its lattice suffix and its suffix phase history are determined by
the successive positions and phases of its relaxed path:
consecutive positions determine the absolute direction. Thus the
map from true suffix histories to relaxed histories is injective even
though raw windows have been quotiented. Summing the comparisons
over histories and lengths up to any finite suffix horizon proves
the finite-horizon assertion. Taking suprema of these nonnegative
finite sums proves the final assertion. No rule is released at a
later time.
\end{proof}

\begin{lemma}\label{l:Bellman-soundness}
Let \(T\) be a nonnegative, finitely branching weighted transition
operator on a set of states, with zero rows on a terminal set.
Let \(g\ge0\) be supported on that terminal set.
If a finite-valued function \(F\ge0\) satisfies
\(g+TF\le F\), then for every integer \(n\ge0\),
\[
 \sum_{j=0}^{n}T^jg\le F .
\]
Consequently the nonnegative sum of all finite first-hit path
weights times their terminal rewards is finite and at most \(F\).
It suffices that \(F=g\) on terminal states and \(TF\le F\)
elsewhere.
\end{lemma}

\begin{proof}
Write \(H_n=\sum_{j=0}^nT^jg\). Each value is a finite sum,
because the operator is finitely branching and \(n\) is finite.
Nonnegativity implies that \(T\) preserves pointwise inequalities.
For \(n=0\), \(H_0=g\le g+TF\le F\). If \(H_n\le F\), then
\[
 H_{n+1}=g+TH_n\le g+TF\le F.
\]
This proves the assertion by induction.

Expanding \(T^jg\) as a finite sum over its successive edges
shows that it counts length-\(j\) paths ending at a state with
terminal reward. A path visiting a terminal state earlier has a
zero outgoing factor, so only first-hit paths have a nonzero
contribution. Conversely each such path appears once, with its
edge-weight product and its terminal reward. The nonnegative sum
over all finite lengths is by definition \(\sup_n H_n\), bounded
by the finite number \(F\) at each initial state. At a terminal
state the row is zero, so \(F=g\) implies the supersolution
inequality there; outside the terminal set \(g=0\), so it is just
\(TF\le F\). This proves the sufficiency statement, including
terminal states of reward zero.
\end{proof}

\begin{definition}\label{d:stationary-certificate}
For a case \((X,d_0)\), fix \(R=70\),
\(B_R=\{-R,\ldots,R\}^2\), \(V=2^{16}\), and \(D_3=2^{48}\).
Supply integers \(Z_{\varphi,q}>0\) and
\(U_{\varphi,q,x}\ge0\), \(x\in B_R\).
Set \(\widetilde U_{\varphi,q,x}=U_{\varphi,q,x}\) inside \(B_R\)
and \(Z_{\varphi,q}\) outside, and \(F=\widetilde U/V\), \(z=Z/V\).
For each individual transition coefficient \(t\) of the exact
stationary or free operator, use the proved rational ceiling
\(\widehat t=p(t)/D_3\), \(p(t)=\lceil D_3t\rceil\);
zero coefficients may be omitted.
The certificate checks all the following:
\[
 \begin{gathered}
 \sum_{\text{free edges }e\text{ from }(\varphi,q)}
       p(t_e)Z_{\operatorname{dest}(e)}
       \le D_3Z_{\varphi,q},\qquad Z_{{\rm C},q},Z_{{\rm D},q}\ge V,\\
 0\le U_{\varphi,q,x}\le Z_{\varphi,q},\qquad
 U_{\varphi,q,c}=0,\qquad U_{\varphi,q,v}=Vg_\varphi,\\
 \sum_{\text{stationary edges }e\text{ from }(\varphi,q,x)}
       p(t_e)\widetilde U_{\operatorname{dest}(e)}
       \le D_3 U_{\varphi,q,x}
       \quad(x\in B_R\setminus\{c,v\}).
 \end{gathered}
\]
All tests are integer inequalities on the regenerated automaton,
directions, phases and lattice sites, not on a trusted supplied matrix.
\end{definition}

An individual transition here fixes a direction and a next phase;
the response pairs belonging to that phase change are summed before
its coefficient is rounded. The definition uses ceilings only for
the finite check; the operators themselves retain their exact
rational coefficients.

For explicit reconstruction of those coefficients, encode directions
by \(0,1,2,3=\mathrm E,\mathrm N,\mathrm W,\mathrm S\).
For an admitted non-revisit edge with incoming direction \(a\)
and outgoing direction \(d\), put
\(\{f,g\}=\{0,1,2,3\}\setminus\{a+2,d\}\), with addition
modulo four. For distinct directions define
\(H(e,e')=1\) for opposite directions and \(H(e,e')=h_C\)
for perpendicular directions. If \(\delta\) specifies the direction
from the site to \(c\), use \(\delta=4\) when \(c\) is not a
neighbour. For an edge not entering \(c\), its response reward is
\[
 K(a,d,\delta)=
 \begin{cases}
 b'_\tau H(g,d),&\delta=f,\\
 b'_\tau H(f,d),&\delta=g,\\
 a_\tau H(f,g)+b_\tau[H(f,d)+H(g,d)],&\delta\notin\{f,g\},
 \end{cases}
 \qquad \tau=\begin{cases}{\rm s},&d=a,\\{\rm t},&d\ne a.\end{cases}
\]
The change-of-phase coefficient is \(x_\tau K/3\). For the free
operator use \(\delta=4\). Diagonal coefficients and revisit
coefficients are those in Definition~\ref{d:phaseweights};
at known sites apply the stated killing and \(1/5\) rules.
If \(t=A/B\ge0\) with integers \(A\ge0,B>0\), its ceiling is
\[
 p(t)=\left\lfloor\frac{D_3A+B-1}{B}\right\rfloor,
 \qquad D_3A\le Bp(t),
\]
and, when \(A>0\), \(B(p(t)-1)<D_3A\). These integer relations
prove \(t\le p(t)/D_3\), including \(t=0\) with \(p(t)=0\).
In particular the diagonal ceiling numerators, in the order
\(x_{\rm s}/3,x_{\rm t}/3,1/3,1/5\), are
\[
 97632666219444,\quad100347585280092,\quad
 93824992236886,\quad56294995342132 .
\]
All response ceilings are specified by the same rational formula
with the exact constants of Lemma~\ref{l:log-mean-bounds} and the
exact dyadic \(h_C\); a decimal replacement of \(h_C\) is not used.

\begin{lemma}\label{l:stationary-local-global}
Every certificate of Definition~\ref{d:stationary-certificate}
defines a finite-valued \(F\) on the entire lattice with
\(F=z\) outside \(B_R\), \(F\le z\), and
\[
 g+T_{X,d_0}F\le F .
\]
For an inside site this follows from its checked integer row and
coefficient ceilings. For an outside site every successor value
is at most its corresponding \(z\), and the exact stationary row
is dominated by the free row; the checked free inequality
implies \(\mathsf Wz\le z\).
Thus \(F\) bounds every stopped continuation as in
Lemma~\ref{l:Bellman-soundness}.
\end{lemma}

\begin{proof}
There are finitely many phases and window classes and finitely many
stored sites. Their integer values, and each positive integer
\(Z_{\varphi,q}\), are finite, so the extension defines a
nonnegative finite-valued \(F\) at every state. The caps prove
\(F\le z\) everywhere; outside the square equality holds.

Divide a checked free integer row by \(D_3V>0\). Since every
coefficient is at most its ceiling and every \(Z\) is nonnegative,
the result implies
\[
 (\mathsf Wz)_{\varphi,q}
 \le\frac1{D_3V}\sum_{e}p(t_e)Z_{\operatorname{dest}(e)}
 \le Z_{\varphi,q}/V=z_{\varphi,q}.
\]
At an inside nonterminal site the identical argument, with
\(\widetilde U\) at each successor, gives \(TF\le F\).
This includes every edge leaving the square; its successor is
assigned the specified \(Z/V\), with no periodic identification.

We justify also the free-row comparison used outside. Ignoring
positions, a non-special stationary edge has the same diagonal
coefficient as the corresponding free edge. Its response coefficient
is no greater: when one non-outgoing neighbour is the centre the
surviving coefficient \(b'_\tau\) is at most the free coefficient
\(b_\tau\), and the other free pair contributions are nonnegative.
Otherwise the response coefficients agree. Removing an edge into
the centre can only decrease a row. At a known site the diagonal
coefficient is reduced to \(1/5\), which is no larger than any
free diagonal coefficient, all response coefficients are removed,
and additional edges may be killed. At a terminal site all outgoing
coefficients vanish. Thus every stationary row is dominated, after
positions are forgotten, by the free row with the same phase and
window, retaining directions and next phases separately.

For an outside site, each successor value is at most its
corresponding \(z_{\psi,q'}\). The preceding comparison therefore
gives
\[
 (TF)(\varphi,q,x)\le(\mathsf Wz)_{\varphi,q}
                       \le z_{\varphi,q}=F(\varphi,q,x).
\]
Both terminals \(c,v\) are inside \(B_{70}\). Their checked
values give \(F=g\) there, with zero outgoing rows; at all other
states \(g=0\). This proves \(g+TF\le F\) everywhere, and
Lemma~\ref{l:Bellman-soundness} applies. The terminal lower bounds
on \(Z\) are also consistent with the target caps and ensure that
the target reward never exceeds its corresponding outside cap.
\end{proof}

\begin{definition}\label{d:prefix-bound}
For each case \((X,d_0)\), let \(S_{X,d_0}\) be the exact
\(h\)-weighted kernel sum over closed trails of lengths \(1\le n\le L=12\)
with first direction \(d_0\), using their prescribed short-trail orderings.
Let \(\mathcal P_{X,d_0,L}\) be all true length-\(L\) prefixes
avoiding \(c\) and \(v\), with first direction \(d_0\).
For \(P\) in this set, let \(p_\varphi(P)\) be the sum of its
exact phase-history products ending in phase \(\varphi\), starting
with weights one in A and B. Use actual trail revisits, actual dead
neighbours and the type-dependent rewards, with the long-trail
straight-first phase weights.
Write \(x(P)\) for its endpoint and \(q(P)\) for its last-window class.
For the potential of its case define the rational number
\[
 J_{X,d_0}=S_{X,d_0}
       +\sum_{P\in\mathcal P_{X,d_0,12}}\sum_\varphi
                    p_\varphi(P)F(\varphi,q(P),x(P)).
\]
A hit at time \(12\) is in \(S_{X,d_0}\), not in the prefix sum.
\end{definition}

Here is a finite recursive definition of the enumeration and of all
its rational summands. It also fixes precisely what is to be
evaluated in Proposition~\ref{p:stationary-check}. A node at time
\(j\le12\) records its full direction word, current position \(x\),
used unoriented bonds \(E\), and counts \(N(y)\) of earlier
departures from every encountered site. Initially these are
\(j=0,x=u,E=\varnothing,N=0\), with incoming direction
\(u-c=\mathrm W\). Its four phase weights are
\((p_{\rm A},p_{\rm B},p_{\rm C},p_{\rm D})=(1,1,0,0)\).
There are three mutually exclusive terminal cases for a node:
\begin{enumerate}
\item If \(x=v\), add its exact short-trail value to \(S_{X,d_0}\)
and produce no children. Test this case before the length cutoff.
\item If \(j=12\) and \(x\ne v\), add its phase-vector dot product
with \(F(\cdot,q(P),x)\) to the prefix sum, and produce no children.
\item Otherwise, try each of the four directions \(d\). At
\(j=0\) retain only \(d=d_0\). At every time discard a move to
\(c\), a move along a used unoriented bond, and the initial
reverse direction. For every remaining move, increment \(N(x)\),
adjoin its bond to \(E\), move to \(x+d\), and recurse.
\end{enumerate}
All children use copies of the parent's data. At a noninitial node
the reverse of its incoming direction is already excluded by the
used-bond test. The true kind is a revisit exactly when the
preincrement value \(N(x)\) is positive. On a first departure
compute \(\kappa=\kappa_h\) using the predeparture counts: a
neighbour \(y\) is dead exactly when \(y=c\) or \(N(y)\ge2\).
The outgoing neighbour of an allowed true extension is not dead.
Use the actual three-, two-, or zero-pair rule of
Definition~\ref{d:kernel}, without a ceiling.

The phase-vector update for a revisit divides all four components
by five. For a straight first departure it is
\[
 (p_{\rm A},p_{\rm B},p_{\rm C},p_{\rm D})\longmapsto
 \left(\frac{x_{\rm s}p_{\rm A}}3,
       \frac{x_{\rm s}p_{\rm B}}3,
       \frac{x_{\rm s}p_{\rm C}}3,
       \frac{p_{\rm D}+x_{\rm s}\kappa p_{\rm B}}3\right);
\]
for a turning first departure it is
\[
 (p_{\rm A},p_{\rm B},p_{\rm C},p_{\rm D})\longmapsto
 \left(\frac{x_{\rm t}p_{\rm A}}3,
       \frac{p_{\rm B}}3,
       \frac{p_{\rm C}+x_{\rm t}\kappa p_{\rm A}}3,
       \frac{p_{\rm D}}3\right).
\]
These formulas use the old components on their right sides.
For a hit at time \(n\le12\), recover the list of actual first
departures \((x_i,\kappa_h(i))\), sort it by the key in
Definition~\ref{d:kernel-order}, including its path-index tie rule,
and evaluate
\[
 3^{-n}(3/5)^{\operatorname{rev}(w)}
       \sum_k\kappa_h(\sigma_w(k))\prod_{l\le k}x_{\sigma_w(l)}.
\]
The phase weights of a short hit are not substituted for this
short-trail expression.

We record why this recursion is exact. Inductively, every generated
node is a true bond trail avoiding \(c\), with no previous hit of
\(v\), except for a terminal node whose endpoint is its first
hit. The used set and departure counts agree with the word by
their update rules. Conversely, every such word of positive length
has its unique shorter parent, and its last direction passes
exactly the tests above. Induction on length therefore proves
completeness without duplication, including all first hits at
length twelve. At any revisited noninitial site each previous
arrival/departure pair uses two distinct bonds. A third departure
would require six bonds. At \(u\), a third departure would require
five, whereas only three avoid \(c\). Thus no branch can depart
from a twice-departed neighbour; equivalently every allowed outgoing
neighbour is live. These observations justify both the departure
counts and the actual dead-site rule at every node. More explicitly,
a twice-departed noninitial site has used all four of its incident
bonds, so a further arrival there is impossible. A twice-departed
\(u\) has used all three bonds available after exclusion of \(c\).
This proves the asserted outgoing-neighbour property even for a
move that would otherwise be the last step of a branch.

Induction on length using the displayed updates proves that the
phase vector is exactly the sum of phase-history products. Keeping
the full word and counts makes each \(\kappa_h(i)\) its value at
the time of that departure; later changes of a neighbour's status
do not change an earlier reward. The short sorting is an exact
comparison of rationals: for positive rewards compare
\[
 (x_i-1)x_j\kappa_h(j)\quad\hbox{and}\quad
 (x_j-1)x_i\kappa_h(i).
\]
Zero rewards have key \(+\infty\). If equal-key entries are grouped
in an implementation, their value is unchanged by exchanging
them: with preceding product \(P\), the difference between the
two orders is
\[
 P\bigl[x_j\kappa_h(j)(x_i-1)
         -x_i\kappa_h(i)(x_j-1)\bigr]=0.
\]
The product following both entries is also unchanged. This proves
agreement with the path-index tie rule even for a grouped evaluation.

There are finitely many branches, of depth at most twelve.
Every coefficient, every potential value and every summand is
rational. For example a common denominator for all possible phase
edge coefficients is
\[
 \Delta=12139925673538045870080000.
\]
It is divisible by the denominators of \(1/3,1/5,x_{\rm s}/3,
x_{\rm t}/3\) and of every possible
\(x_\tau\kappa_h/3\) from the two types and the four choices of
dead non-outgoing neighbours. This is a finite rational divisibility
check with the displayed constants. Multiplying phase weights by
\(\Delta^j\) makes every update integral. If \(A\) is the sum
of the resulting length-twelve phase numerators times the
corresponding stored \(U\) entries, its exact tail contribution is
\(A/(\Delta^{12}V)\). Adding the independently evaluated rational
short part specifies \(J_{X,d_0}\) by finite integer arithmetic.

For a compact endpoint lookup, every length-twelve endpoint satisfies
\(|x-u|_1\le12\) and
\((x_1-u_1)+(x_2-u_2)\equiv12\pmod2\): each nearest-neighbour
step changes this parity and changes distance by at most one.
All these sites lie in \(B_{70}\). A compact table must copy the
entry with spatial index
\((x_1+70)141+(x_2+70)\) from the full potential, for every phase
and referenced window class. This identity, checked for every
referenced coordinate, justifies the lookup. Parity is used only
to identify possible endpoints, and introduces no numerical factor.

\begin{lemma}\label{l:prefix-domination}
For every case whose integer certificate passes, the \(h\)-weighted
contribution of all true closed trails of that first direction,
over every length, is at most \(J_{X,d_0}\).
This holds uniformly for every finite length cutoff, with the same
upper bound. The short part uses the exact order through \(12\);
all longer trails split uniquely into a surviving length-\(12\)
prefix and a dominated stopped suffix.
\end{lemma}

\begin{proof}
Fix a finite total-length cutoff \(R_0\). Trails of length at most
\(\min(R_0,12)\) contribute a sub-sum of \(S_{X,d_0}\), because
their weights are nonnegative and their orderings are exactly the
specified short orderings. This includes any hit at time twelve.

If \(R_0>12\), every remaining true trail has a unique first twelve
steps \(P\), which avoid \(c\) and \(v\), and a suffix of positive
length at most \(R_0-12\), stopping at its first hit of \(v\).
Its prescribed ordering is straight-first, so
Lemma~\ref{l:phase-expansion} applies to its full contribution.
Split each phase history after those twelve departures and group
by its phase \(\varphi\). The sum of the prefix factors is exactly
\(p_\varphi(P)\). Lemma~\ref{l:stationary-domination} bounds the
sum of its true suffix factors and terminal rewards by the
corresponding stationary stopped sum from
\((\varphi,q(P),x(P))\). This comparison is valid for the chosen
cutoff and with the full true prefix as initial history.
Lemmas~\ref{l:stationary-local-global} and
\ref{l:Bellman-soundness} bound that stopped sum by
\(F(\varphi,q(P),x(P))\), independently of the cutoff.

There are finitely many length-twelve prefixes. For each, multiplying
by its nonnegative phase weights and then summing yields the prefix
term in \(J_{X,d_0}\). Even if two prefixes have the same class,
endpoint and phase, their separate prefix weights remain in the
sum, so quotienting loses no multiplicity. Adding the short bound
gives the assertion for \(R_0\). For \(R_0\le12\), including
\(R_0=0\), the nonnegative prefix term may simply be added to the
short bound. Taking the supremum over all finite cutoffs proves
the all-length assertion.
\end{proof}

\begin{lemma}\label{l:kernel-reflection}
For row B, reflection \((x_1,x_2)\mapsto(x_1,-x_2)\) fixes
\(c,u,v\), interchanges the first directions N and S, and preserves
every trail weight, type, dead pair, shift, \(h\)-reward and
prescribed ordering. Hence the true B/N and B/S kernel contributions
are equal at every length and every finite cutoff.
No corresponding reduction is asserted for row C.
\end{lemma}

\begin{proof}
Let \(r(x_1,x_2)=(x_1,-x_2)\). In row B it fixes the three
specified sites and the incoming direction \(\mathrm W\), and
interchanges \(\mathrm N\) and \(\mathrm S\). It maps each
unoriented nearest-neighbour bond bijectively to such a bond. Thus
\(w\mapsto rw\) is a bijection, inverse to itself, between the
two first-direction classes, preserving length, centre avoidance,
first-hit stopping and the set of path indices of earlier
departures. At every index it preserves straightness, turning,
revisit status and every departure count.

It maps the three non-arrival neighbours to the corresponding
neighbours of the reflected departure, and fixes the centre as a
set. Hence a neighbour is dead if and only if its image is dead.
It preserves the outgoing neighbour and may interchange the two
other slots, whose rules are symmetric. The actual assigned
coefficients \(a_\tau,b_\tau,b'_\tau\) therefore agree for each
corresponding live pair. Opposite pairs remain opposite and
perpendicular pairs remain perpendicular, so their types and
\(h\)-weights agree. Since \(r\) is an isometry and fixes \(u,v\),
\[
 \max_{z\in p}\min\{d(z,u),d(z,v)\}
 =\max_{z\in rp}\min\{d(z,u),d(z,v)\}.
\]
Thus shifts agree as well.

Each first departure has the same rational ordering key on both
trails. Its path index is unchanged, so the short ordering including
tie breaking is preserved. The long straight-first ordering is
preserved for the same reason. All products and response sums
in Definition~\ref{d:kernel} consequently agree. Summing the
bijection at a fixed length, and then over any finite length cutoff,
proves the assertion. No equality of automaton classes or of
stationary potentials under \(r\) is needed.
\end{proof}

\begin{proposition}\label{p:stationary-check}
For the five cases B/N, B/W, C/N, C/W, C/S, there exist
integer arrays \(Z,U\) satisfying every check in
Definition~\ref{d:stationary-certificate}, with
\(m=10,L=12,R=70\) and the specified scales.
Their regenerated quotient has \(3741\) states and the spatial
square has \(19881\) sites. Their exact rational prefix evaluations obey
\[
 \begin{array}{c|ccccc}
 (X,d_0)&(B,{\rm N})&(B,{\rm W})&(C,{\rm N})&(C,{\rm W})&(C,{\rm S})\\ \hline
 10^9J_{X,d_0}\ \le&
 351207554&290599163&527872005&380754375&273048088
 \end{array}
\]
Each displayed upper bound denotes an exact integer inequality
after clearing the rational prefix denominator.
\end{proposition}

\begin{proof}[Finite data and evaluation contract]
The five witnesses are the files
\[
 \texttt{stationary\_}\!X\texttt{\_m10\_R70\_d}j
                  \texttt{\_coarse.bin.gz},
 \qquad (X,j)=(B,1),(B,2),(C,1),(C,2),(C,3),
\]
in the archived data directory. Here \(1,2,3\) mean
\(\mathrm N,\mathrm W,\mathrm S\). After decompression each file
has nine little-endian signed 64-bit header integers
\[
 [10,12,70,3741,19881,v_1,v_2,16,j],
\]
where \((v_1,v_2)=(1,0)\) for B and \((0,1)\) for C.
It then has \(4\cdot3741\) unsigned 32-bit integers for \(Z\),
followed by \(4\cdot3741\cdot19881\) unsigned 32-bit integers
for \(U\), and no other payload. Phases are ordered A,B,C,D,
then class numbers in the order of
Proposition~\ref{p:automaton10}; for \(U\), the last index is
\((x_1+70)141+(x_2+70)\). This identifies each supplied literal
with a mathematical array index. In particular,
\(19881=141^2\), rather than a periodic spatial state count.

The finite assertions on these arrays are exactly every integer
inequality and equality in Definition~\ref{d:stationary-certificate}.
For each of the \(4\cdot3741\) phase/class indices, require
\(Z>0\), its free row inequality, and the C,D lower terminal
bounds. For each of the \(4\cdot3741\cdot19881\) stored
phase/class/site indices, require \(0\le U\le Z\), the specified
centre and target equalities, and the nonterminal stationary row
inequality when appropriate. The coefficients are regenerated
from the exact formulas following that definition, not inferred
from floating values stored in a program.

For clarity, all four direction choices, both diagonal and
change-of-phase possibilities, and every admitted successor are
included in those rows. At a known site remove the initial bond
directions and every detected-revisit direction, use \(1/5\) on
the remaining same-phase edges, and include no response edge.
At the centre and target check the terminal values instead of an
outgoing row. For an outside successor substitute its \(Z\) entry.
The quotient transitions and representative incoming directions
are justified by the raw reconstruction in
Proposition~\ref{p:automaton10}. In particular a class with mixed
incoming directions can be used only on revisit transitions, for
which no response direction is needed.

There are \(4\cdot3741\cdot(19881-2)=297469356\) nonterminal
local component inequalities per file, as well as all the free,
cap and terminal checks just specified. These are exact integer
inequalities over the five supplied arrays; gzip checksums or the
termination status of a native checker are not substitutes for them.

Evaluate the finite recursion following
Definition~\ref{d:prefix-bound} against these same arrays, retaining
their exact class and coordinate indices. Completeness of that
recursion and the phase-vector update have already been proved
there. The resulting finite enumeration checks, in the column
order of the proposition, are
\[
 \begin{array}{c|rrrrr}
 &B/{\rm N}&B/{\rm W}&C/{\rm N}&C/{\rm W}&C/{\rm S}\\ \hline
 \text{nodes}&139985&161929&105025&155673&145485\\
 \text{surviving prefixes}&89150&103257&66947&99168&92748\\
 \text{short first hits}&937&922&715&1036&861.
 \end{array}
\]
Counts alone are not the prefix assertion: the check must also
evaluate every rational summand. If the exact short sum is
\(s/d\), \(d>0\), and the nonnegative integer prefix accumulator
is \(A\) with common denominator \(\Delta^{12}V\) as specified
there, the displayed bound \(10^9J\le K\) means exactly
\[
 10^9\bigl(s\Delta^{12}V+Ad\bigr)
                  \le Kd\Delta^{12}V .
\]
Use respectively
\(K=351207554,290599163,527872005,380754375,273048088\).
This gives the finite rational prefix tests without a decimal
rounding convention. A compact lookup is accepted only with the
entry-by-entry identity with the full potential described after
Definition~\ref{d:prefix-bound}.
The archived exact fractions in the five prefix records identify
the values of these finite expressions; the expressions and the
cleared inequalities above, together with all the array checks,
are the numerical assertions to be verified separately.
\end{proof}

\begin{lemma}\label{l:stationary-assembly}
With the five certified potentials and their exact prefix bounds,
\[
 (M_1h)_B\le2J_{B,{\rm N}}+J_{B,{\rm W}}\le\rho_0h_B,\qquad
 (M_1h)_C\le J_{C,{\rm N}}+J_{C,{\rm W}}+J_{C,{\rm S}}
                                      \le\rho_0h_C ,
\]
where \(h_C=5404319552844595/4503599627370496\) and
\(\rho_0=993015/10^6<1\).
The same inequalities hold for every finite true-trail cutoff;
therefore (H3\(_h\)) holds for the modified kernel.
\end{lemma}

\begin{proof}
Every true trail from \(u\) has first direction N, W or S: direction
E enters the forbidden centre. For row B,
Lemma~\ref{l:kernel-reflection} identifies the true N and S sums
at every finite cutoff; apply Lemma~\ref{l:prefix-domination}
to B/N and B/W. For row C apply that lemma separately to its
three first-direction classes, with no symmetry reduction.
This proves the first inequalities in the statement, uniformly in
the total-length cutoff.

The exact component ceilings of Proposition~\ref{p:stationary-check}
give
\[
 2J_{B,{\rm N}}+J_{B,{\rm W}}
 \le\frac{2\cdot351207554+290599163}{10^9}
 =\frac{993014271}{10^9}
 <\frac{993015000}{10^9}=\rho_0h_B.
\]
For C their sum is
\[
 J_{C,{\rm N}}+J_{C,{\rm W}}+J_{C,{\rm S}}
 \le\frac{1181674468}{10^9}
 <\frac{119}{100}
 <\frac{993015}{10^6}
             \frac{5404319552844595}{4503599627370496}.
\]
The middle inequality is \(1181674468<1190000000\).
For the last inequality, use the exact rational comparisons
\[
 h_C>1199/1000,\qquad
 993015\cdot1199=1190624985>1190000000.
\]
The first comparison is an exact comparison with the displayed
dyadic fraction. These two inequalities imply
\(\rho_0h_C>119/100\). All comparisons involve positive integer
denominators. In particular they give the desired C-row bound.
Finally \(0<\rho_0<1\) since \(0<993015<1000000\), and
\(h_B,h_C>0\).

The same upper bounds hold for every finite cutoff. By
Definition~\ref{h:H3}, this proves (H3\(_h\)). Equivalently, taking
suprema of the nonnegative row sums gives a finite \(M_1h\)
with the asserted bound. Only a finite prefix was used in the
numerical evaluation; the Bellman inequality controls all suffix
horizons simultaneously.
\end{proof}

